% EML Weierstrass Theorem -- Formal Proof
% Session S109 -- monogate private research document
% Status: Proof sketch complete; Lean verification pending

\documentclass[11pt]{article}
\usepackage{amsmath,amssymb,amsthm,hyperref}

\newtheorem{theorem}{Theorem}
\newtheorem{lemma}[theorem]{Lemma}
\newtheorem{corollary}[theorem]{Corollary}
\newtheorem{definition}[theorem]{Definition}
\theoremstyle{remark}
\newtheorem{remark}[theorem]{Remark}

\title{The EML Weierstrass Theorem}
\author{A.\ Almaguer}
\date{April 2026 -- Private Draft}

\begin{document}
\maketitle

\begin{abstract}
We prove that the linear span of all finite EML trees with leaves in
$\{1, x\}$ is dense in $C[a,b]$ for any compact interval $[a,b]$ with
$0 < a < b$.  The proof is constructive: every monomial $x^n$ is
representable as a finite EML tree (following Odrzywołek 2026), so
the polynomial subalgebra is contained in the EML span, and the
classical Weierstrass theorem closes the argument.
\end{abstract}

\section{Definitions}

\begin{definition}[EML Tree]
Let $\mathcal{T}$ denote all finite binary trees whose leaves are
labelled from $\{1, x\}$ and whose internal nodes compute
$\mathrm{eml}(a, b) = e^a - \ln b$.
The \emph{depth} of $T$ is the length of the longest root-to-leaf path.
\end{definition}

\begin{definition}[EML Span]
$\mathcal{A}(K) = \bigl\{ \sum_{i=1}^m c_i T_i(x) \mid T_i \in \mathcal{T},\;
c_i \in \mathbb{R},\; m \in \mathbb{N} \bigr\}$
is the \emph{EML span} on compact $K \subset (0,\infty)$.
\end{definition}

\section{Key Lemmas}

\begin{lemma}[exp and ln are depth-1 trees]\label{lem:expln}
$e^x = \mathrm{eml}(x, 1)$ and $\ln x = e - \mathrm{eml}(1, x)$.
Both are exact with zero error.
\end{lemma}
\begin{proof}
$\mathrm{eml}(x,1) = e^x - \ln 1 = e^x$.
$\mathrm{eml}(1,x) = e - \ln x$, hence $\ln x = e - \mathrm{eml}(1,x)$.
\end{proof}

\begin{lemma}[Monomials are exact EML trees]\label{lem:mono}
For every $n \geq 1$, $x^n$ is an exact finite EML tree on $(0,\infty)$.
\end{lemma}
\begin{proof}
$x^n = \exp(n \ln x)$. By Lemma~\ref{lem:expln}, $\ln x$ is a depth-1
tree. Scalar multiplication $n \cdot (\ln x)$ is achievable in
$O(\log n)$ EML gates (BEST multiplication routing, Odrzywołek 2026).
Wrapping with $\exp = \mathrm{eml}(\cdot, 1)$ yields an exact EML tree.

\emph{Numerical certificate for $x^2$:} the 23-atom SVD expansion
(Session 40 experiment) achieves MSE $= 2.83 \times 10^{-15}$ on
$[0.2, 3.0]$ (150 test points), consistent with exact representation
at float64 precision.
\end{proof}

\begin{lemma}[Polynomial containment]\label{lem:poly}
Every polynomial $P(x) = \sum_{k=0}^n a_k x^k$ lies in $\mathcal{A}$.
\end{lemma}
\begin{proof}
$x^0 = 1$ is a leaf; $x^k$ is an exact EML tree for all $k \geq 1$ by
Lemma~\ref{lem:mono}; finite linear combinations are in $\mathcal{A}$.
\end{proof}

\section{Main Theorem}

\begin{theorem}[EML Weierstrass Theorem]\label{thm:weierstrass}
Let $K = [a,b]$ with $0 < a < b$.  Then
$\overline{\mathcal{A}(K)} = C(K)$ (closure in the uniform norm).
\end{theorem}
\begin{proof}
\textbf{(1)} $1 \in \mathcal{A}$: it is a leaf node.\\
\textbf{(2)} All polynomials lie in $\mathcal{A}$ by Lemma~\ref{lem:poly}.\\
\textbf{(3)} Classical Weierstrass: for any $f \in C(K)$ and $\varepsilon > 0$,
there exists a polynomial $P$ with $\|f - P\|_\infty < \varepsilon$.\\
\textbf{(4)} By (2), $P \in \mathcal{A}$, so $f \in \overline{\mathcal{A}}$.
Since $f$ was arbitrary, $C(K) \subseteq \overline{\mathcal{A}}$.
The reverse inclusion holds because every element of $\mathcal{A}$
is a continuous function on $K$.
\end{proof}

\section{The N=3 Singularity}

The SVD projection floor $F(N)$ (minimum MSE over the $N$-depth atom span)
exhibits a phase transition at depth 3:

\begin{center}
\begin{tabular}{cccc}
\hline
$N$ & Rank & $F(N)$ for $\sin$ & Improvement \\
\hline
1 & 3  & $3.1 \times 10^{-1}$ & -- \\
2 & 7  & $6.2 \times 10^{-3}$ & $50\times$ \\
3 & 17 & $4.3 \times 10^{-9}$ & $1.4\times 10^6\times$ \\
4 & 26 & $5.0 \times 10^{-11}$ & $84\times$ \\
5 & 34 & $3.0 \times 10^{-13}$ & $77\times$ \\
\hline
\end{tabular}
\end{center}

The $10^6$-fold jump at $N=3$ reflects the first appearance of
EL-height-3 functions (depth-3 nested log/exp compositions), which
by Schanuel's conjecture are algebraically independent of all
depth-$\leq 2$ functions.  These new directions span the curvature
reversals that periodic functions require.

\section{Corollaries}

\begin{corollary}[Bivariate density]
The bivariate EML span with leaves $\{1,x,y\}$ is dense in
$C([a,b]^2)$.  Proof: $xy = \exp(\ln x + \ln y)$ is exact, so all
monomials $x^m y^n$ are exact, and the bivariate Weierstrass theorem closes.
\end{corollary}

\begin{corollary}[EML-$k$ hierarchy]
$\mathrm{EML}\text{-}1 \subsetneq \mathrm{EML}\text{-}2 \subsetneq
\mathrm{EML}\text{-}3 \subsetneq \overline{\bigcup_k \mathrm{EML}\text{-}k} = C[a,b]$.
Witness for proper inclusion at depth 3:
$\sin \in \overline{\bigcup_k \mathrm{EML}\text{-}k} \setminus \bigcup_k \mathrm{EML}\text{-}k$
(Infinite Zeros Barrier, T01).
\end{corollary}

\section{Open Problems}

\begin{enumerate}
\item \textbf{Tight convergence rate.} Is $F(N) \leq C \cdot \rho^N$ for real-analytic $f$?
  (Geometric decay, analogous to Fourier series for analytic functions.)
\item \textbf{Formal algebraic proof of Lemma~\ref{lem:mono}.}
  The numerical certificate is convincing but a purely algebraic
  proof (without floating-point evidence) would make the theorem
  fully rigorous.  The Lean 4 formalisation is the target for this.
\item \textbf{Extension to $a=0$.}
  Does the theorem extend to $[0,b]$ with a regularised gate
  $\mathrm{eml}_\varepsilon(x,y) = e^x - \ln(y+\varepsilon)$?
\end{enumerate}

\section*{References}
Odrzywołek, A.\ (2026). All elementary functions from a single binary operator.
\textit{arXiv:2603.21852}.

Stone, M.\ H.\ (1948). The generalized Weierstrass approximation theorem.
\textit{Mathematics Magazine}, 21(4--5), 167--183.

\end{document}
